\section{产物的结构与组成表征}
\label{sec:validation-endpoints}

图~\ref{fig:route-matrix}右侧列出的四类表征分别回答不同问题，不能互相替代：单晶或先进衍射确认结构模型，PXRD/Rietveld判断批量相组成，化学分析和显微观察检查实际元素比与污染，高温实验界定特定气氛和热史下的稳定范围。表~\ref{tab:synthesis-conditions}的“产物与表征”栏据此限定每条实验记录可以支持的结论。缺少某类数据时，只对相应结论保持谨慎，不能用其他方法补足。

\subsection{结构模型}

单晶X射线衍射可把目标相判断落实到实测晶体的结构模型；目前的直接研究正是据此鉴定$\mathrm{Ba_5Y_{12}Zn[O(SiO_4)]_8}$\cite{ababaikeri2024ba5y12zn}。在无Zn相关Ba--Y硅酸盐中，单晶精修揭示了非整比调制结构，CRED、同步辐射PXRD和中子精修的组合确认了独立的四方Ba--Y相区\cite{yamane2024synthesis,gulay2024navigation}。其他单晶或同步辐射研究确立了各自Ba--Y、Ba--Zn--Si化合物的结构，但结论只适用于相应化合物\cite{motozawa2022bay16si4o33,zou2021crystal}。单晶结构能证明被测晶体与精修模型一致，却不能单独证明整批样品相纯，也不能把名义投料组成等同于实测占位；局域谱学至多补充多型辨析，不提供目标相的直接成相证据\cite{becerro2004revisiting}。

\subsection{批量相组成}

PXRD/Rietveld面向批量样品，重点分析相分数、未索引峰及相同热史下的晶型差异。Ba--Y陶瓷研究通过粉末衍射发现次生相，独立相区研究把同步辐射粉末数据与其他衍射方法结合确认相组成\cite{yamane2024microstructure,gulay2024navigation}；Ba--Zn--Si研究也利用X射线、中子或同步辐射数据区分批量晶型和结构模型\cite{lin1999phase,zou2021crystal}。因此，PXRD/Rietveld可支持批量相组成和杂相判断，但峰位相似或出现主相都不能自动推出相纯度，也不能替代单晶占位判断。对只报告固相转变或热致变化的研究，本文只借鉴其表征方法，不把题名中的“优化”当作目标相判据\cite{kerstan2012thermal,cai2024optimized}。

\subsection{实际组成与污染}

名义配比、局部测点和样品整体组成需要分别记录。EDS/EDX、EPMA或ICP可检查晶粒间成分差异、外来相和处理介质带入。无ZnBa--Y陶瓷研究把次生相与玛瑙来源的$\mathrm{SiO_2}$污染同时纳入分析\cite{yamane2024microstructure}；独立Ba--Y相区研究以EDS配合多种衍射方法约束组成和结构\cite{gulay2024navigation}；Ba/Sr--Zn--Si相关化合物则用单晶结构和EDX联合确认自身组成关系\cite{thieme2015ba1}。这些方法只能说明测点或取样体积内检测到的元素，不能凭少量测点判断整批样品均匀性，也不能把烧前名义配比直接改写成烧后化学式。显微、热学和光谱联测可为掺杂或晶体生长提供过程参照，但不承担目标相组成结论\cite{tejas2024structural,pang2005study}。

\subsection{高温稳定性}

HT-XRD、热分析、膨胀测量和多气氛实验回答材料在特定热史或气氛下如何变化，不是室温结构鉴定的替代。BaZn$_2$Si$_2$O$_7$的低温和高温晶型由中子与X射线衍射区分，说明原位高温相不能等同于冷却后产物\cite{lin1999phase}；Zn/Co及Ba/Sr--Zn--Si/Ge相关研究进一步把取代组成与热稳定性或相边界联系起来\cite{kerstan2013bazn2si2o7,thieme2016negative}。相关陶瓷在特定还原气氛中的耐受性只能作气氛参照，不能证明目标相具有相同稳定性\cite{zou2019anti}。高温实验可界定已发表体系中的转变、分解或气氛耐受范围，但原位信号不能单独证明冷却后样品相纯；缺少相应数据时，稳定性结论保持未确定。

四类表征共同限定结论的适用范围：衍射确认结构模型，粉末衍射判断批量相组成，化学分析检验实际元素比和污染，高温实验界定热史与气氛边界。只有把这些结果与同一批次的合成条件对应起来，才能讨论目标相是否复现，以及相关工艺可以外推到何种范围。
